% strategic_diff/drafts/zh/manuscript.tex
% P4.2 中文 manuscript 骨架 (plan-r3 衍生论文扩展, 2026-05-15)
% 编译: xelatex + biber + xelatex × 2

% papers_assets/_shared/preamble_zh.tex
% Shared Chinese LaTeX preamble (XeLaTeX + CJKutf8).
% P4.1 (plan-r3 derivative-paper extension, 2026-05-15)
% Compile: xelatex + biber + xelatex × 2

\documentclass[12pt,a4paper]{article}

% --- CJK Fonts (XeLaTeX, cross-platform Q1.1b fallback) ---
\usepackage{fontspec}
\usepackage{xeCJK}
% Windows (local dev): SimSun/SimHei/FangSong
% Linux/macOS (CI): Noto Serif CJK SC / Noto Sans CJK SC (fonts-noto-cjk apt package)
% Last resort: AR PL UMing CN (fonts-arphic-uming) which is preinstalled in many texlive distros
\IfFontExistsTF{SimSun}{%
  \setCJKmainfont{SimSun}[BoldFont=SimHei]%
}{%
  \IfFontExistsTF{Noto Serif CJK SC}{%
    \setCJKmainfont{Noto Serif CJK SC}[BoldFont=Noto Serif CJK SC]%
  }{%
    \setCJKmainfont{AR PL UMing CN}[AutoFakeBold=2]%
  }%
}
\IfFontExistsTF{SimHei}{%
  \setCJKsansfont{SimHei}%
}{%
  \IfFontExistsTF{Noto Sans CJK SC}{%
    \setCJKsansfont{Noto Sans CJK SC}%
  }{%
    \setCJKsansfont{AR PL UKai CN}%
  }%
}
\IfFontExistsTF{FangSong}{%
  \setCJKmonofont{FangSong}%
}{%
  \IfFontExistsTF{Noto Sans Mono CJK SC}{%
    \setCJKmonofont{Noto Sans Mono CJK SC}%
  }{%
    \setCJKmonofont{AR PL UMing CN}%
  }%
}
\setmainfont{Times New Roman}[Ligatures=TeX]
\xeCJKsetup{CJKecglue={\hskip 0.15em plus 0.05em minus 0.05em}}

% --- Math ---
\usepackage{amsmath, amssymb, amsthm}

% --- Tables & Figures ---
\usepackage{booktabs}
\usepackage{graphicx}
\usepackage{longtable}
\usepackage{rotating}
\usepackage{caption}
\captionsetup{font=small,labelfont=bf,skip=4pt}

% --- Layout ---
\usepackage[margin=1in]{geometry}
\usepackage{setspace}
\onehalfspacing
\usepackage{hyperref}
\hypersetup{colorlinks=true,linkcolor=blue,citecolor=blue,urlcolor=blue}

% --- Bibliography ---
\usepackage[backend=biber,style=apa,natbib=true]{biblatex}

% --- Macros ---
\newcommand{\code}[1]{\texttt{#1}}
\usepackage[dvipsnames]{xcolor}
\newcommand{\NoteThis}[1]{\textbf{\textcolor{red}{[TODO: #1]}}}

% --- Theorem environments ---
\newtheorem{prop}{命题}
\newtheorem{hyp}{假设}


\addbibresource{refs.bib}

\title{银行金融科技报道中的战略差异度：\\多方法文本测度及其实证表现}
\author{\NoteThis{作者姓名}}
\date{\today}

\begin{document}
\maketitle

\begin{abstract}
\NoteThis{最终摘要待定; 结果性语句以生产面板回归为准。}
金融科技冲击下, 中国银行业的数字化转型叙事呈现高度同质化, 差异化战略能否
转化为特许权价值成为悬而未决的实证问题。本文将
Hoberg--Phillips (2016) 内生产品差异化框架移植到中国银行业金融科技竞争场景,
基于银行—年度金融科技媒体报道语料, 采用 Guzman--Li (2023) 的 N-最近邻
文本距离并辅以产品市场流动性 (Hoberg--Phillips--Prabhala 2014) 等多种文本
表征, 测度各银行战略叙事的差异度 (\code{strat\_diff\_n5} 为预设主测度)。
识别上, 我们采用银行与年份双向固定效应、银行层面聚类标准误的
单一预设验证性 (confirmatory) 主规格, 辅以 Bartik 型移位—份额工具变量与
Romano--Wolf 族错误率校正的推断纪律。理论预期为: 战略差异度提升上市银行的
特许权价值 (Tobin Q, $\beta>0$), 并经由存款定价能力压低资金成本
($\beta<0$)。我们发现 \NoteThis{核心结论 (生产面板 T=19 回归结果待填)}\@。
\end{abstract}

\section{引言}\label{sec:intro}

过去十余年, 金融科技从支付、信贷到财富管理全方位重塑了银行业的竞争格局
\parencite{berg2022fintechlending,parlour2022payments,vives2025lendercompetition},
中国是这一变革最剧烈的市场之一 \parencite{hau2024fintechcredit,cheng2020bankfintech}。
面对科技平台的进入压力, 中国商业银行几乎同步宣示"数字化转型"战略,
零售转型、开放银行、数字中台等叙事在同业间高度趋同。当所有银行都讲述
相似的故事时, 一个经典的产业组织问题重新浮现: 战略差异化是否仍然能够
削弱同质竞争、积累市场势力, 并最终体现为特许权价值溢价
\parencite{hoberg2016tnic,drechsler2026franchise}?
媒体报道是观察银行战略叙事的天然窗口: 新闻文本记录了各银行金融科技
布局的方向、强度与措辞, 且媒体覆盖本身即是资本市场的重要信息中介
\parencite{fang2009media,you2018pen}。本文构建覆盖 2007--2025 年的中国银行
金融科技新闻语料: 原始普查 (过滤、去重之前的全量) 计 18{,}198{,}978 条
媒体报道、来自 3{,}740 家媒体来源, 其中报纸类报道占全部行的 92.93\,\%
并被保留; 银行—金融科技实体共现硬过滤前, 生产目标语料取普查
总量口径 18{,}198{,}978 条 (2007--2025), 覆盖 110 家银行、
19 个年度 (T=19); 共现硬过滤为丢行式硬过滤, 过滤后进入回归面板的
文章数远小于该目标语料量。

本文的研究问题是: 银行金融科技战略叙事的\emph{文本差异度}
(distinctness) 是否提升其\emph{特许权价值}?
我们将 \textcite{hoberg2016tnic} 的内生产品差异化框架移植到银行业
金融科技竞争场景, 借鉴 \textcite{guzman2023strategy} 的 N-最近邻文本距离
方法, 在银行—年度层面测度每家银行的战略叙事与其 N 个最相似同业之间的
平均距离 (\code{strat\_diff\_n5}, N=5、TF-IDF 字符表征为预设主测度),
并以产品市场流动性 \parencite{hoberg2014fluidity}、行业质心距离、
LDA 主题散度与嵌入距离等多种表征互为稳健性。理论逻辑是:
差异化程度高的银行面临的同质竞争更弱、产品市场威胁更低,
从而积累市场势力, 在上市银行子样本中体现为 Tobin Q 溢价;
在全样本通道上, 差异化经由存款定价能力压低资金成本。

识别上, 我们施加银行与年份双向固定效应, 在银行层面聚类标准误,
并预先声明单一验证性主规格 (\code{strat\_diff\_n5} $\to$ Tobin Q,
单侧理论预期 $\beta>0$), 将全部替代测度、替代轨道与异质性分析
明确标注为探索性 (exploratory), 族内以 Romano--Wolf 逐步降步法
控制族错误率 \parencite{romano2005stepwise,list2019multiplehypothesis}。
针对媒体口径与企业价值之间的经典内生性担忧 \parencite{gurun2012hype},
我们进一步引入 Bartik 型移位—份额工具变量, 其构造与诊断遵循份额外生
\parencite{goldsmith2020bartik} 与冲击外生
\parencite{borusyak2022shiftshare,borusyak2025shiftshareguide} 两条路径的
现代标准。

本文有四方面贡献。其一, 测度贡献: 据我们所知, 本文首次在中文银行新闻
语料上落地 N-最近邻战略差异度测度, 并给出五种文本表征
(TF-IDF 字符 / TF-IDF 分词 / LDA+JSD / 嵌入 / 行业质心) 乘
N-邻梯子 (N$\in\{3,\dots,7,10,20,30,40,50\}$) 的完整测度族,
延伸了文本战略测度谱系
\parencite{guzman2023strategy,bellstam2021innovation,banker2024strategy,hoberg2025scope}。
其二, 识别与推断纪律贡献: 观测性银行面板研究中显式划定
confirmatory/exploratory 边界、预设族错误率校正并配 Bartik 工具的
做法仍不多见。其三, 实证贡献: 为金融科技时代银行差异化战略与
特许权价值的关系提供来自中国的系统证据, 补充存款特许权文献
\parencite{drechsler2026franchise} 的竞争战略维度。其四, 方法论贡献:
针对文本测度结论对嵌入模型选择的敏感性 \parencite{ganguli2024patenttext},
本文所有嵌入派生测度均在两个独立厂商的 co-primary 嵌入模型下
并行计算、互为稳健性 (方法节 F5)。

本文隶属于同一数据管线上的系列论文, 但研究对象彼此正交:
与语调论文 (P1) 不同, 本文测度的是 N-邻文本差异度而非公司语调;
与关注度论文 (P2) 不同, 本文测度差异度而非曝光强度; 本文不做
实质—象征二元标签 (P3), 不区分创新区位 (P4); 与 benefit\_society
的单一态度投影不同, 本文使用多文本表征差异; 与 brokerage\_network
的网络位置测度不同, 本文使用文本表征差异; 与 llm\_peer 的 LLM
语义同业工具不同, 本文的核心是 N-邻文本距离。

下文结构如下: 第 \ref{sec:lit} 节综述相关文献; 第 \ref{sec:theory}
节发展假设; 第 \ref{sec:method} 节介绍数据管线与测度构造;
第 \ref{sec:data} 节描述样本与变量; 第 \ref{sec:main} 节报告基准
回归; 第 \ref{sec:mech} 节检验机制; 第 \ref{sec:het} 节分析异质性;
第 \ref{sec:robust} 节呈现稳健性; 第 \ref{sec:concl} 节总结。

\section{文献综述}\label{sec:lit}

\subsection{文本战略差异度与最优区分度}
本文测度构念的直接源头是两支文献。其一是 \textcite{hoberg2016tnic}
开创的文本网络行业 (TNIC) 谱系: 以企业文本的成对相似度刻画内生
产品差异化, 差异化程度决定企业面临的竞争强度;
\textcite{hoberg2014fluidity} 进一步以竞争者文本向邻域移动的速度
定义产品市场流动性 (fluidity), 度量企业面临的产品市场威胁;
\textcite{hoberg2025scope} 将该谱系延伸到 21 世纪企业的范围与规模。
其二是 \textcite{guzman2023strategy} 的创业战略测度: 以企业与其 N 个
最近邻之间的文本距离度量战略差异度 (distinctness), 本文的
\code{strat\_diff\_n\{N\}} 列族即其银行新闻版实现。相关地,
\textcite{bellstam2021innovation} 用分析师报告文本测度企业创新,
\textcite{banker2024strategy} 用文本法刻画战略类型与成本结构。
最优区分度 (optimal distinctiveness) 文献则提供边界条件:
\textcite{haans2019distinctiveness} 表明差异化的绩效回报依类别同质性
而变; \textcite{su2024optimal} 强调企业需在竞争压力与制度压力之间
动态定位; \textcite{arts2025technology} 给出技术差异化缓解产品市场
竞争的最新证据。这些边界条件对应本文的所有制、规模与网络枢纽
异质性设计。行业质心距离法在管线代码内部另有一条中文文献线索
(谢佳松等, 2025, 《经济研究》), 该题录尚待 CNKI 核实, 故本文正式
引注以 \textcite{hoberg2016tnic} 与 \textcite{guzman2023strategy} 锚定
质心与 N-邻两类操作化。

\subsection{金融科技、银行竞争与特许权价值}
金融科技进入如何重塑银行竞争是近年金融学的核心议题:
\textcite{berg2022fintechlending} 综述了金融科技信贷的兴起;
\textcite{vives2025lendercompetition} 从理论上刻画信息技术对贷款人
竞争的影响; \textcite{parlour2022payments} 分析金融科技争夺支付流
对银行信息租金的侵蚀。中国情境下, \textcite{hau2024fintechcredit}
提供金融科技信贷促进创业增长的证据; \textcite{cheng2020bankfintech}
基于中文语料构造银行金融科技指数并考察其信用风险后果, 是与本文
文本测度做法最直接的对标。outcome 端, 本文的"差异化 $\to$ 市场势力
$\to$ 特许权价值"叙事与存款特许权文献相承:
\textcite{drechsler2026franchise} 表明存款特许权是银行价值的核心
组成部分, 其稳定性依赖银行维持存款定价能力——这正是本文
资金成本通道的理论落点。

\subsection{Text-as-data 方法与嵌入模型选择敏感性}
\textcite{gentzkow2019text} 给出"文本 $\to$ 数值表示 $\to$ 回归输入"
的方法总纲。\textcite{hanley2019dynamic} 与 \textcite{engle2020climatenews}
分别示范了在金融申报文本与新闻文本上构造时变暴露测度的范式;
\textcite{fang2009media} 确立媒体覆盖测度与资产价格研究的经典设计,
\textcite{you2018pen} 则刻画了中国国有与市场化媒体的制度分野,
为本文使用中文媒体语料提供情境正当性。嵌入时代的前沿方法上,
\textcite{cong2024textualfactors} 提出可扩展、可解释的文本因子框架;
尤为关键的是, \textcite{ganguli2024patenttext} 系统论证了文本相似度
结论强烈依赖嵌入模型的选择——这一发现直接支撑本文
Qwen3-Embedding-8B 与 text-embedding-3-large 双主轨严格列对偶设计
(方法节 F5) 的学理必要性。

\subsection{识别策略与推断纪律}
本文的工具变量策略立足移位—份额 (Bartik) 识别文献:
\textcite{goldsmith2020bartik} 给出份额外生路径及 Rotemberg 权重诊断,
\textcite{borusyak2022shiftshare} 与 \textcite{borusyak2025shiftshareguide}
给出冲击外生路径与实操指南; 二者共同支撑本文 Bartik 工具三轨各自
恰好识别的设定与诊断报告。多重检验纪律上,
\textcite{romano2005stepwise} 的逐步降步法是本文族错误率校正的
方法源, \textcite{list2019multiplehypothesis} 提供实验与观测研究中
confirmatory/exploratory 边界声明的操作先例。媒体内生性方面,
\textcite{gurun2012hype} 关于地方媒体倾向与企业价值的经典警示
正当化了本文的工具变量需求; 聚类推断遵循
\textcite{cameron2015practitioner}。

综合来看, 既有文献尚未在银行业——尤其是金融科技竞争最激烈的中国
市场——落地过 N-最近邻文本战略差异度测度, 也缺少将该测度与
特许权价值联系起来、并施加现代推断纪律的系统证据。本文填补这一
缺口。

\section{理论与假设}\label{sec:theory}

\subsection{差异化与特许权价值}
在 \textcite{hoberg2016tnic} 的内生产品差异化框架中, 企业文本与
竞争者越相似, 其面临的替代威胁越强、定价能力越弱; 反之, 差异化
使企业占据更稀疏的产品空间位置, 削弱同质竞争并积累市场势力。
移植到银行业: 当多数银行以高度趋同的叙事推进金融科技转型时,
战略叙事显著区别于同业的银行更可能占据差异化的业务生态位
(客群、场景、技术路线), 其市场势力沉淀为特许权价值, 在资本市场
定价中体现为估值溢价。特许权价值在上市银行中可由 Tobin Q 度量。
据此提出理论预期 (而非结果):

\textbf{H1 (confirmatory).} 战略差异度提升银行特许权价值:
\code{strat\_diff\_n5} 对 Tobin Q (回归实现为缩尾后的
\code{tobin\_q\_w}) 的系数预期为正 ($\beta>0$); 上市银行子样本主表。
此为本文唯一预设验证性主规格。

\subsection{差异化与资金成本}
特许权价值的微观基础之一是存款定价能力
\parencite{drechsler2026franchise}: 面临同质竞争弱的银行无需以更高
利率争夺存款。若差异化确实削弱同质竞争, 其应在负债端体现为更低的
资金成本。该通道可在全样本 (含非上市银行) 中检验, 且其方向预测与
本管线语调论文 (P1) 的语调—资金成本通道相反, 二者互为交叉验证。

\textbf{H2.} 战略差异度降低资金成本: \code{strat\_diff\_n5} 对
\code{funding\_cost} 的系数预期为负 ($\beta<0$); 全样本通道。

\subsection{收入结构与产品市场威胁通道}
差异化若有效, 应可观测到两条中间通道。其一, 差异化银行更可能在
非利息业务 (财富管理、科技输出、场景金融) 上形成独特定位, 收入
结构向非利息收入倾斜。其二, 按产品市场流动性逻辑
\parencite{hoberg2014fluidity}, 差异化银行所处产品空间中竞争者向其
移动的速度更慢, 面临的产品市场威胁更低。盈利能力 (ROA) 为辅助
通道。

\textbf{H3 (机制, 探索性).} \code{strat\_diff\_n5} 对
\code{noninterest\_income\_share} 预期为正; 对 \code{strat\_fluidity}
预期为负; 对 \code{roa} 预期为正。均为 secondary outcome, 探索性
标注。

\subsection{差异化与网络枢纽地位的互补性}
差异化战略的价值可能依赖银行在同业网络中的位置: 既差异化又居于
枢纽位置的银行, 兼具独特生态位与信息、业务网络优势, 构成最强
护城河。但网络位置与差异化的交乘项内生, 无法赋予因果解释
\parencite{brambor2006understanding}。

\textbf{H4 (exploratory, 条件关联).} \code{diff\_x\_btwn}
(= \code{strat\_diff\_n5} $\times$ L.\code{bank\_between\_cent}) 与
Tobin Q 预期呈正向\emph{条件关联}; 按审计裁决
(STRATDIFF-MOD-1), 该交乘仅作条件关联陈述, 绝不作因果表述。

\section{方法}\label{sec:method}
\subsection{样本与数据}
银行—年度面板覆盖约 164 个银行—年观测 (2017--2018)。
所有数据由 \code{news\_process} 管线 (Phase 9b 基线) 处理,
详见 \code{papers\_assets/V6\_METHODOLOGY\_CHANGES.md}。

\subsection{语调标注 (F1: 三家族 LLM 真三角验证)}
采用三个独立 LLM 家族: \code{gpt-5.4-nano} (OpenAI)、
\code{gemini-3.1-flash-lite} (Google)、\code{deepseek-v4-flash}
(DeepSeek)。OpenAI 与 Google 两家通过 OpenLux
(\code{/v1/chat/completions}, OpenAI 兼容) 调用, DeepSeek 直连
\code{api.deepseek.com}\footnote{平台迁移说明: 本文数据的 LLM 推理于 2026 年 7 月经本管线
当时的聚合中转代理完成; 自 2026 年 8 月起, 管线的聚合中转统一迁移至
OpenLux (\code{https://api.openlux.ai/v1})。缓存键由「模型名 + prompt
文本 + prompt 版本 + temperature」四元组构成, 不含端点地址, 故本次迁移
对既有推理结果逐位无影响, 复现包内的确定性缓存可原样离线重放。}; 所有调用 \code{temperature=0}、\code{seed=42},
并使用字节级冻结 prompt (版本 v8)。
家族间一致性以 Cohen's $\kappa$ 与 Krippendorff's $\alpha$ 报告;
多数投票得最终标签; 平局/空票触发最多 80 轮确定性重试, 仍未破局则
保守编码为\code{中性} (而非遗留为 undefined)。

\subsection{篇章级编码 (F2: 去污染)}
三家族全部失败者编为\emph{缺失} (非 \code{neutral}); 三票分裂/空票
不计入缺失, 触发最多 80 轮确定性重试仍未破局则保守编码为\code{中性};
银行--金融科技相关性未通过多数验证者亦视为缺失。
此举消除了 v6 之前 81.7\,\% ``中性'' 标签的污染。

留篇门控仅取全文轨 (\code{llm\_verified\_full}); 实体上下文轨
(\code{ctx}) 与银行—金科同句共现轨 (\code{ss}) 为以全文轨为条件
(conditional-on-full) 的稳健性轨, 其行宇宙是全文轨留篇集的子集。
相较早期"任一轨验证通过即保留"的版本, 更严格的门控会机械性压缩
\code{n\_articles} 及派生份额 (\code{fintech\_share}、
\code{attention\_pct}) --- 这是定义性、方向偏保守的水平位移,
并非数据错误。

\subsection{零观测语义 (F3: NaN 主路径)}
对 \code{n\_articles=0} 的银行—年, 相似度类特征
(\code{auth\_sim\_*}、\code{benefit\_*}、\code{article\_sim\_*})、
比例类 (\code{token\_adopt\_*}、\code{hidden\_adopt\_*}、
\code{subst\_adopt\_*})、\code{fintech\_intensity\_mean} 均编为
\code{NaN}; 计数类保持 0。回归采用列删除 (Allison 2001 / Heckman 1979)。

\subsection{Bartik IV 修剪 (F4)}
约 30 个候选 Bartik 工具变量中, 28 个在 v6 中因违反
Goldsmith-Pinkham \emph{et~al.}\ (2020, \emph{AER}) 外生性被删除。
保留: \code{bartik\_type\_iv\_full} 与 \code{prov\_leave1out\_tone\_full}。

\subsection{嵌入类测度 (双主轨、互为稳健性; F5)}
所有嵌入派生特征均在来自独立厂商的两个同为主测度的嵌入模型下并行
计算、并排报告、互为稳健性: 阿里 \code{Qwen3-Embedding-8B}
(4{,}096 维, 经硅基流动 API; canonical 无后缀列) 与 OpenAI
\code{text-embedding-3-large} (3{,}072 维; \code{\_lg} 后缀列),
遵循严格列对偶 (任一特征两轨齐备或皆无)。各轨进入各自恰好识别的
独立设定, 绝不合并进同一过度识别系统。早期管线版本 canonical 槽为
Google Gemini 嵌入; 自管线 v3.15.0 (2026 年 6 月) 起换为阿里
\code{Qwen3-Embedding-8B}, 全部嵌入派生特征经受控全量重嵌再生,
嵌入空间语义效度以 Spearman 秩相关闸把关 ($\rho \geq 0.80$)。
情感标注 LLM 集成 (F1) 不受此变更影响: 嵌入模型与情感 LLM
是互不相交的管线组件。

\subsection{战略差异度测度构造}\label{sec:method-stratdiff}
本文专属的处理变量族 \code{strat\_diff\_*} 在银行—年度层面按五种
文本表征构造, 各表征下的差异度定义一致: 给定某年各银行的金融科技
新闻文本聚合向量, 银行 $i$ 的战略差异度为其与 $N$ 个最近邻同业之间
的平均距离 (N-最近邻生态位距离, 参照 \textcite{guzman2023strategy}
的操作化)。五种表征为: (i) TF-IDF 字符 n-gram (canonical,
预设主测度 \code{strat\_diff\_n5} 即 N=5 的字符表征);
(ii) TF-IDF 分词 (jieba); (iii) LDA 主题分布, 距离取
Jensen--Shannon 散度; (iv) 嵌入向量余弦距离 (双主轨, 见 F5);
(v) 行业质心法: 银行文本向量与留一 (leave-one-out) 行业质心的
余弦距离 (\code{centroid\_diff\_*} 列族)。辅助测度包括与最近邻的
最大相似度 \code{strat\_sim\_max}、银行内文本一致性
\code{coh\_tfidf\_c\_full}, 以及产品市场流动性 \code{strat\_fluidity}
——按 \textcite{hoberg2014fluidity} 的流动性构造逻辑, 以竞争者文本
向本行邻域的移动刻画产品市场威胁。

\subsection{N-邻梯子与 FWER 族定义}\label{sec:method-fwer}
邻域宽度 N 无先验唯一值, 本文将 N$\in\{3,4,5,6,7,10,20,30,40,50\}$
全部物化 (\code{strat\_diff\_n\{N\}}), 各 N 互为稳健性、不分主次,
各自进入恰好识别的独立规格。多重检验按预设族校正: headline
n\_neighbor 族 = 字符表征全部 10 个 N 成员; centroid 族 = 4 个
质心变体 (\code{centroid\_diff\_tfidf\_c}、\code{centroid\_diff\_tfidf\_w}、
\code{centroid\_diff\_lda}、\code{centroid\_diff\_emb}); 族内以
Romano--Wolf 逐步降步重抽样 \parencite{romano2005stepwise} 报告
校正后 $p$ 值。word/lda/emb/\code{\_lg} 及 ctx/ss 轨的大 N 列仅作
点估计探索性稳健, 不组建独立 FWER 族。

\subsection{聚类与推断口径}\label{sec:method-cluster}
主规格标准误在银行层面单向聚类 (生产面板 $G\approx110$,
即 110 家银行、T=19); 上市银行子样本主表的聚类数更少, 以
Cameron--Gelbach--Miller 类小 $G$ 敏感性检验兜底; 稳健性检验
对比银行 $\times$ 年份双向聚类
\parencite{cameron2015practitioner}。confirmatory 主规格仅有一条
(H1), 其余全部替代测度、轨道与子样本均标注 exploratory。

\subsection{网络控制列与缺失值约定}\label{sec:method-network}
为区分"文本战略差异"与"网络孤立"两个混淆构念, 主规格控制银行的
Burt 结构约束 (\code{bank\_constraint}, 网络嵌入度)。披露约定:
度数 $<2$ 的银行—年该列为 \code{NaN}, 回归列删除带来的样本损失
须在表注中报告 (N-02); 其余中心性列对脱网银行按 zero-fill 约定
填零并披露 N。网络共现判据 (同篇 / $\pm$1 句窗 \code{\_win1} /
严格同句 \code{\_sent}) 的三档变体用于控制列替换检验
(见第 \ref{sec:robust} 节)。

\subsection{研究者—AI 协作披露}\label{sec:method-ai}
遵循目标期刊 AI 使用政策, 作者披露三类 AI 工具的使用方式, 全部
输出均由作者审核确认、文责自负。

\emph{LLM 情感标注与路由。}文章级情感与相关性标注采用三家独立
厂商 LLM 的多数投票集成 (OpenAI \code{gpt-5.4-nano}、Google
\code{gemini-3.1-flash-lite}、DeepSeek \code{deepseek-v4-flash});
OpenAI 与 Google 两家经 OpenLux 的 OpenAI 兼容端点路由, DeepSeek
直连 \code{api.deepseek.com}。所有调用 \code{temperature=0}、
\code{seed=42}、prompt 版本 v8 字节级冻结; 完整调用日志与缓存
数据库保存在复制包中, 支持零成本离线复现。

\emph{留篇门控。}一篇文章是否计为银行—金融科技报道仅由全文轨
验证 (\code{llm\_verified\_full}) 决定; \code{ctx} 与 \code{ss} 轨为
conditional-on-full 的稳健性轨 (详见 F2)。该口径的机械性水平位移
(压缩 \code{n\_articles} 及派生份额) 系定义性后果, 已在 F2 披露。

\emph{嵌入双轨与厂商换模。}全部嵌入派生测度在阿里
\code{Qwen3-Embedding-8B} (4{,}096 维, canonical) 与 OpenAI
\code{text-embedding-3-large} (3{,}072 维, \code{\_lg}) 两个
co-primary 模型下并行计算、严格列对偶、互为稳健性; canonical 槽
的厂商维度经历 Google $\to$ 阿里 Qwen 换模 (管线 v3.15.0, 2026 年
6 月), 换模后全部嵌入派生特征受控全量重嵌再生并通过 Spearman
$\rho\geq0.80$ 语义效度闸 (详见 F5)。

\section{数据描述}\label{sec:data}

\subsection{语料口径链条}
本文语料经历三级口径, 须严格区分。第一级为原始普查口径 (过滤、
去重之前的全量): 183 个去重源文件、18{,}198{,}978 条媒体报道,
时间跨度 2007-01-01 至 2025-12-31, 覆盖 3{,}740 家媒体来源;
媒体类型为报纸的报道占全部行的 92.93\,\% (16{,}912{,}361 行) 被
保留, 通讯社、杂志、政府公报与网页等其余类型按行剔除。第二级为
生产目标语料口径: 取普查总量、银行—金融科技
实体共现硬过滤\emph{前}, 18{,}198{,}978 条 (2007--2025), 覆盖
110 家银行、T=19。第三级为回归面板口径: 经共现硬过滤 (丢行式)
与留篇门控后的银行—年度面板, 过滤后文章条数远小于第二级目标语料
\NoteThis{共现硬过滤+留篇门控后的文章条数待生产机全量重跑后实读
填写。}; 本稿开发环境下为约 164 个银行—年观测
(2017--2018, 早期试跑 T=2 开发切片), 生产面板以 110 家银行 $\times$ 19 年
为目标。\NoteThis{生产面板最终 N 待财务面板 enrichment 接入与
全量重跑后填写。}

\subsection{变量定义}
\emph{处理变量}: \code{strat\_diff\_n\{3,4,5,6,7,10,20,30,40,50\}}
$\times$ 五方法 (TF-IDF 字符 / TF-IDF 分词 / LDA+JSD / 嵌入 /
质心) $\times$ 三文本轨 (full/ctx/ss), 嵌入方法另有 \code{\_lg}
孪生列; 预设主测度为 \code{strat\_diff\_n5} (N=5, TF-IDF 字符,
full 轨)。\emph{结果变量}: 主结果为上市银行子样本的
\code{tobin\_q} (回归实现为缩尾 \code{tobin\_q\_w}); 次级结果为
\code{funding\_cost}、\code{roa}、\code{strat\_fluidity} 与
\code{noninterest\_income\_share}。\emph{控制变量}: 银行层面 5 列
(\code{ln\_assets}、\code{leverage}、\code{capital\_ratio}、
\code{npl\_ratio}、\code{deposit\_growth}) 与年度宏观 2 列
(\code{gdp\_growth}、\code{m2\_growth}); 财务面板 enrichment 即将
接入, 当前复现包控制集为新闻覆盖代理 (\code{ln\_news\_coverage}、
\code{media\_hhi}、\code{bank\_scale\_ord})。网络控制为
\code{bank\_constraint} (见第 \ref{sec:method-network} 节)。

\subsection{描述统计}
\NoteThis{Table 1 描述统计参见 \code{papers/strategic_diff/tables.py:build\_table1};
均值/标准差/分位数等全部数值待生产面板重跑后填入, 不得先行填写。}

\section{基准回归}\label{sec:main}

验证性主规格为上市银行子样本上的双向固定效应回归:
\begin{equation}\label{eq:main}
\mathit{TobinQ}^{w}_{it}
  = \beta\,\mathit{StratDiff}^{(n5)}_{it}
  + \boldsymbol{\gamma}'\mathbf{X}_{it}
  + \alpha_i + \lambda_t + \varepsilon_{it},
\end{equation}
其中 $\mathit{TobinQ}^{w}_{it}$ 为缩尾后的 Tobin Q,
$\mathit{StratDiff}^{(n5)}_{it}$ 为 N=5 TF-IDF 字符表征的战略差异度,
$\mathbf{X}_{it}$ 为银行层面与年度宏观控制, $\alpha_i$、$\lambda_t$
分别为银行与年份固定效应; 标准误在银行层面聚类。Stata 实现为
\code{reghdfe tobin\_q\_w strat\_diff\_n5 \$controls,
absorb(bank\_id year) vce(cluster bank\_id)}。按 H1 的单侧理论预期
$\beta>0$。Table 2 将按 M1--M6 渐进控制列报告式 (\ref{eq:main}),
从无控制、仅固定效应, 逐列加入覆盖代理、财务控制与网络控制;
Figure 2 为对应系数图。若 $\beta$ 在加入全部控制后保持正向且
经济量级可观, 则支持"差异化 $\to$ 特许权价值"命题。
全样本通道以 \code{funding\_cost} 为被解释变量重估式
(\ref{eq:main}), 按 H2 预期 $\beta<0$。
\NoteThis{Table 2 全部系数、标准误与显著性待生产面板 (T=19)
回归结果填入; 源 \code{output/tables/table2\_*.csv}。}

针对媒体口径内生性 \parencite{gurun2012hype}, 我们以 Bartik 型
移位—份额工具作两阶段最小二乘:
\begin{equation}\label{eq:iv}
\widehat{\mathit{Text}}_{it}
  = \pi\,Z^{\mathrm{bartik}}_{it}
  + \boldsymbol{\phi}'\mathbf{X}_{it} + \alpha_i + \lambda_t + u_{it},
\end{equation}
第一阶段报告 Kleibergen--Paap rk F 统计量, 并按
\textcite{goldsmith2020bartik} 报告 Rotemberg 权重诊断
(Table A3)。式 \eqref{eq:iv} 中 $\widehat{\mathit{Text}}_{it}$
即本文的内生处理变量——主测度 \code{strat\_diff\_n5} (full 轨);
工具与内生变量沿用管线 IV 扩展的同轨配对约定, 每条文本轨的内生
变量配同轨工具 (\code{bartik\_type\_iv\_full} 为 primary), 各轨
各自恰好识别; 严禁跨轨合成过度识别系统 (详见第 \ref{sec:robust}
节)。\NoteThis{\code{bartik\_type\_iv\_full} 在管线 IV 扩展登记中
原为 full 轨语调测度构造的工具; 其对 \code{strat\_diff\_n5} 的
第一阶段强度须以生产面板 KP rk F 实证确认后方可作为本文 IV 主
规格。IV 估计值与第一阶段 F 待生产面板结果填入。}

\section{机制检验}\label{sec:mech}

若差异化经市场势力起效, 应观测到两条通道。其一为存款定价通道:
差异化银行面临的同质竞争弱, 无需以高利率争夺存款, 资金成本更低
(H2); 该通道方向与本管线语调论文 (P1) 的语调通道相反, 二者互为
交叉验证——语调捕捉媒体评价的态度维度, 差异度捕捉战略定位的
空间维度, 两者若在资金成本上给出相反且各自符合理论的符号,
可相互排除"媒体覆盖量的机械效应"这一共同混淆解释。其二为产品
市场威胁与收入结构通道 (H3): Table 3 将分别以
\code{strat\_fluidity} (预期负)、\code{noninterest\_income\_share}
(预期正) 与 \code{roa} (预期正) 为被解释变量重估式
(\ref{eq:main}), 均标注探索性。解读上, 若 fluidity 显著为负而
非利息收入份额显著为正, 则支持"差异化占据稀疏生态位、收入结构
独特化"的机制链条。
\NoteThis{Table 3 (源 \code{output/tables/table3\_*.csv}) 与
Figure 3 数值待生产面板结果填入。}

\section{异质性}\label{sec:het}

按最优区分度文献的边界条件
\parencite{haans2019distinctiveness,su2024optimal}, Table 4 与
Table 8 将报告四类子样本拆分: 国有 (SOE) vs.\ 非国有、网络枢纽
(hub) vs.\ 非枢纽、规模分组与银行类型分组; 全部标注 exploratory。
理论上, 制度压力更强的国有大行差异化空间受限, 差异化回报可能
集中于中小股份制与城商行; 该预期属探索性, 不作预设假设检验。

调节效应方面, 按 H4 估计
\begin{equation}\label{eq:mod}
\mathit{TobinQ}^{w}_{it}
  = \beta_1\,\mathit{StratDiff}^{(n5)}_{it}
  + \beta_2\,\mathit{StratDiff}^{(n5)}_{it}\times
    L.\mathit{BankBetween}_{i,t-1}
  + \beta_3\,L.\mathit{BankBetween}_{i,t-1}
  + \boldsymbol{\gamma}'\mathbf{X}_{it} + \alpha_i + \lambda_t
  + \varepsilon_{it},
\end{equation}
交互项解读遵循 \textcite{brambor2006understanding} (报告边际效应
而非单看交乘系数)。按审计裁决 (STRATDIFF-MOD-1), 网络位置与
差异化的交乘内生, 式 (\ref{eq:mod}) 的 $\beta_2$ 仅作\emph{条件
关联}陈述——"既差异化又居枢纽的银行—年伴随更高估值"——绝不
表述为因果。控制列 \code{bank\_constraint} 作为嵌入度混淆控制
进入全部异质性规格; 其 deg$<2\to$NaN 的列删除样本损失在表注
披露 (N-02)。
\NoteThis{Table 4 / Table 8 数值待生产面板结果填入。}

\section{稳健性}\label{sec:robust}

本节各稳健性维度共享同一推断纪律: 每一梯子内的各成员\emph{互为
稳健性、不分主次}, 各自进入恰好识别的独立规格并排报告;
任何梯子的成员都\emph{绝不}合并进同一过度识别系统 (如以多个
高相关孪生测度或多条工具同时进入一个 Hansen J 检验)——高相关
孪生并非独立信息, 合成过度识别检验只会制造虚假的过度识别功效。

\paragraph{N-邻敏感性梯子 (headline FWER 族)。}
Table 5a 以 \code{strat\_diff\_n\{3,4,5,6,7,10,20,30,40,50\}}
(TF-IDF 字符, full 轨) 逐一单跑式 (\ref{eq:main}), 族内报告
Romano--Wolf 校正后 $p$ 值 \parencite{romano2005stepwise}。两点
统计披露: 其一, 2026-07 扩展将族成员由 5 (n3--n7) 增至 10
(纳入 n10--n50), Romano--Wolf 校正随族扩大\emph{单调变严},
既有 n3--n7 (含 n5) 的校正后 $p$ 值因此被机械抬高; 这是扩族的
定义性后果, 审稿人应据校正后口径解读。其二, \code{strat\_diff\_n5}
具有双角色: 它既是 confirmatory 主规格中\emph{未经校正的单一预设
假设} (第 \ref{sec:main} 节), 又是本族内接受 RW 校正的普通成员。
两个口径正交、服务不同目的——前者检验单一预设理论假设, 后者
回答"N 的选择是否驱动结论"的族问题——同一测度出现在两处并非
重复计数 (double-counting)。word/lda/emb/\code{\_lg} 及 ctx/ss 轨
的大 N 列仅作点估计探索稳健, 不组建独立 FWER 族。

\paragraph{表征方法族。}
五种文本表征 (TF-IDF 字符 / TF-IDF 分词 / LDA+JSD / 嵌入 / 质心)
各自单跑、绝不同时进入同一右端; headline 保持字符表征。质心族
(\code{centroid\_diff\_tfidf\_c} / \code{centroid\_diff\_tfidf\_w} /
\code{centroid\_diff\_lda} / \code{centroid\_diff\_emb}) 为第二个预设
FWER 族 (4 成员), 同样报告 RW 校正后 $p$。辅助测度
\code{strat\_sim\_max} 与 \code{coh\_tfidf\_c\_full} 各自单跑。

\paragraph{三文本轨 full/ctx/ss。}
留篇门控为 full-only, \code{ctx} 与 \code{ss} 轨是
conditional-on-full 的稳健性轨, 其行宇宙为 full 轨留篇集的子集
(F2)。\code{strat\_diff\_n5\_ctx} 与 \code{strat\_diff\_n5\_ss}
各自单跑式 (\ref{eq:main}); 三轨互为稳健性、高相关孪生非独立
假设, 严禁 full$\times$ctx$\times$ss 合成 Hansen J。

\paragraph{co-primary 双嵌入。}
嵌入方法列族在 Qwen3-Embedding-8B (canonical) 与
text-embedding-3-large (\code{\_lg}) 下严格列对偶并排 (F5);
两轨各进恰好识别设定、禁混同一过度识别系统。学理必要性:
\textcite{ganguli2024patenttext} 表明文本相似度结论强烈依赖嵌入
模型选择, 跨厂商双轨若给出一致符号, 结论即对该选择稳健。

\paragraph{Bartik IV 三轨镜像与诊断。}
IV 规格按轨镜像: \code{bartik\_type\_iv\_full} (primary) /
\code{bartik\_type\_iv\_ctx} (fallback) / \code{bartik\_type\_iv\_ss}
(ss 镜像), 每轨内生变量配同轨工具、各自恰好识别; 严禁机械
\code{\_ctx}$\to$\code{\_ss} 替换合成过度识别。约 30 个候选工具
已按 \textcite{goldsmith2020bartik} 外生性标准修剪至 2 (F4)。
Table A3 报告 Rotemberg 权重与 KP rk F。
\NoteThis{各轨 IV 点估计与诊断量待生产面板填入。}

\paragraph{IV 三引擎并排。}
除 Bartik 外, 管线另提供同组留一均值工具
(\code{leave\_one\_out\_peer\_*\_iv}) 与媒体来源 house-style 暴露
工具 (\code{outlet\_*\_iv}); 三类工具各自单独 2SLS、互为识别
策略稳健性, 绝不同时进入同一方程
\parencite{borusyak2025shiftshareguide}。

\paragraph{聚类推断。}
主规格银行单向聚类; Table 5 / Table 7 对比银行 $\times$ 年份
双向聚类 \parencite{cameron2015practitioner}。

\paragraph{网络控制列判据梯子。}
网络共现判据三档 (同篇, 无后缀 / $\pm$1 句窗 \code{\_win1} /
严格同句 \code{\_sent}) 互为稳健性、各自恰好识别并排 (规格
稳健性并排, 不设 FWER 族)。本文用到的 \code{bank\_constraint}
混淆控制与 \code{diff\_x\_btwn} 调节项均有三判据孪生, 稳健性
检验以逐一替换控制列/交乘列重估主规格; 同时披露
\code{bank\_constraint} 的 deg$<2\to$NaN 列删除损样本 (N-02)
与其余中心性列的 zero-fill 约定。

\paragraph{分时代稳定性。}
T=19 生产面板允许按金融科技兴起前后分时代重估主规格
(Table 9, 探索性), 检验主效应的时变稳定性。
\NoteThis{分时代切点与结果待生产面板确定与填入。}

\paragraph{管线级梯子的披露边界。}
情感 LLM 三家、BERT 三模型、政治词权重、词典消融等管线级稳健性
梯子作用于情感/词典列族, 与本文的文本距离处理变量无直接交集,
故不进入本文稳健性 battery, 仅在方法节 (F1) 作管线披露。

\section{结论}\label{sec:concl}

本文将内生产品差异化框架移植到中国银行业金融科技竞争场景,
在中文银行新闻语料上首次落地 N-最近邻文本战略差异度测度,
并以单一预设验证性规格、Bartik 工具与 Romano--Wolf 族错误率
校正的推断纪律, 检验"差异化 $\to$ 市场势力 $\to$ 特许权价值"
命题。我们发现 \NoteThis{核心结论回顾 (生产面板结果待填)}\@。

理论上, 本文把最优区分度与文本产业组织两支文献带入银行业情境:
若差异化确实伴随估值溢价与更低资金成本, 则说明在强监管、高同质
的银行业, 战略叙事的差异化仍是特许权价值的可行来源。政策上,
结果 \NoteThis{视符号与量级而定} 将提示监管者: 引导银行差异化
定位 (而非同质化竞速式的数字化军备竞赛) 有助于维护银行
特许权价值与体系稳定; 对银行管理层, 差异化的生态位选择是
金融科技时代的战略杠杆。

本文有三点局限。其一, 面板为年度频率, 战略叙事的年内动态与
短期市场反应无法捕捉。其二, 媒体报道本身内生: 报道口径受媒体
选择与银行公关影响 \parencite{gurun2012hype}; 尽管我们以 Bartik
与留一工具缓解, 工具的排他性约束不可直接检验, 因果解释仍需
审慎。其三, 样本限于中国单一市场, 制度环境 (国有主导、强监管、
平台经济格局) 特殊, 外推到其他银行体系须谨慎。财务面板
enrichment 与全量 T=19 生产回归即将接入, 届时更新全部
结果性陈述。

\NoteThis{最终讨论与政策含义细化待结果落定后补写。}

\printbibliography

\end{document}
